% SPDX-License-Identifier: CC-BY-SA-4.0
% Author: Matthieu Perrin
% Part: <Nom de la partie>
% Section: <Nom de la section>
% Sub-section: <Nom de la sous-section>  % (facultatif, laisser vide si non utilisé)
% Frame: <Titre de la slide>

\begingroup

\begin{frame}{De la calculabilité à la complexité}

  \on[width=1.05\textwidth, x=-4mm, y=27mm]{
    \myquote{Scott Aaronson}{
      La question de savoir si $\textsc{p} \stackrel{?}{=} \textsc{np}$ est l'analogue moderne, en temps fini,\\
      de la question de la décidabilité.
      \vspace{1mm}\linebreak
      Dans les deux cas, on demande : \structure{``Peut-on automatiser la créativité ?''}
    }
  }

  \onImage[y=3mm, x=-43mm]{%
    width=2cm,
    title={Juris Hartmanis\footnote{Prix Turing 1993 avec Richard Stearns pour leurs travaux sur la complexité}},
    licenselogo={\ccbysa},
    license={\href{https://creativecommons.org/licenses/by-sa/3.0/}{CC BY-SA 3.0} -- Wolfgang Hunscher 2002 -- \href{https://commons.wikimedia.org/wiki/File:Juris_Hartmanis(2002).jpg}{Wikimedia}},
    img={Hartmanis.jpg}
  }

  
  \on[width=.7\textwidth, x=9mm, y=1mm]{
    \myquote{Juris Hartmanis}{
      Les structures que nous avons trouvées \linebreak
      dans \textsc{np}
      sont \structure{les ombres exactes} \linebreak
      des structures découvertes par Post et Kleene dans les degrés de Turing
    }
  }


  \on[width=1.1\textwidth, x=-4mm, y=-27mm]{
    \myquote{Christos Papadimitriou}{
      Un langage est dans \textsc{re} s'il existe un certificat pour son appartenance.\linebreak
      Il est dans \textsc{np} si ce certificat est polynomial.\vspace{1mm}\linebreak
      La différence n'est pas dans la nature de la preuve,\linebreak
      mais dans \structure{la patience du vérificateur}.
    }
  }

\end{frame}

\endgroup

